\section{The $\tau$ factors and Lemma H}
\label{app:tau}

\begin{lemma}[Lemma H] \proofst
Split $h-1=e^{-w^2/2s^2}(R_{\rm lo}+R_{\rm hi})$ at frequency $m$.
For the kernels $K=2\sum\kappa_n\cos nt$ of App.~D of~\cite{F},
\begin{equation}
|\langle K(h-1)\rangle|\le\eta_{\rm lo}\|K\|_2+\eta_{\rm hi}\|K_{\ge m}\|_2 .
\end{equation}
Since $\kappa_n\propto n^ke^{-n^2\tau/2+n\beta}$, the high part carries $e^{-m^2\tau/2}$, and $m^2\tau=u^2/b\ge u^2$.
\end{lemma}

\emph{$\tau$ cells.} \caproofst{\code{cb\_tau.py}}
The $\tau$-dependent factors of the bound decrease in $\tau$ and increase in $\beta$.
On $400$ cells of $[20,80]$, with $\beta\le10^{-6}\tau$ and an audit of $20000$ points, their product is at least $0.99977$.
On $[40,80]$ it is at least $1-1.1\times10^{-8}$.
In certificate (ii) of Sec.~\ref{sec:cert} the $\nu$'s are used directly at $\tau=20$.
This is valid on the whole range $\tau\ge20$ \proofst.
Indeed $\tau^j\nu_k(\tau,10^{-6}\tau)=2\sum_n n^k\tau^je^{-c_n\tau}$ with $c_n=n^2/2-10^{-6}n\ge1/2-10^{-6}$, and
\begin{equation}
\frac{d}{d\tau}\,\tau^je^{-c\tau}=\tau^{j-1}e^{-c\tau}\,(j-c\tau)<0\quad\text{for }\tau>j/c .
\end{equation}
Since $j\le2$, $j/c_n<4.0001<20$, so every term decreases on $\tau\ge20$.
For $\omega_2=2\sum_ne^{-n^2\tau/2}(n^2\tau-1)$ (here $n^2\tau\ge20$) each term has the derivative $(n^2/2)e^{-n^2\tau/2}(3-n^2\tau)<0$.
A $60$-digit evaluation on $4000$ points of $[20,10^8]$ finds no increase \numst; with $\beta=0.6\tau$ in place of $10^{-6}\tau$ (so that $c_1<0$) it finds an increase at every point (\code{cq\_checks.py}).
